\documentclass[12pt,a4paper,twoside,openright]{report}
\let\openright=\cleardoublepage



%%% Choose a language %%%

\newif\ifEN
% \ENtrue   % uncomment this for english
\ENfalse   % uncomment this for czech

%%% Configuration of the title page %%%

\def\ThesisTitleStyle{mff} % MFF style
%\def\ThesisTitleStyle{cuni} % uncomment for old-style with cuni.cz logo
%\def\ThesisTitleStyle{natur} % uncomment for nature faculty logo

\def\UKFaculty{Faculty of Mathematics and Physics}
%\def\UKFaculty{Faculty of Science}

\def\UKName{Charles University in Prague} % this is not used in the "mff" style

% Thesis type names, as used in several places in the title
\def\ThesisTypeTitle{\ifEN BACHELOR THESIS \else BAKALÁŘSKÁ PRÁCE \fi}
%\def\ThesisTypeTitle{\ifEN MASTER THESIS \else DIPLOMOVÁ PRÁCE \fi}
%\def\ThesisTypeTitle{\ifEN RIGOROUS THESIS \else RIGORÓZNÍ PRÁCE \fi}
%\def\ThesisTypeTitle{\ifEN DOCTORAL THESIS \else DISERTAČNÍ PRÁCE \fi}
\def\ThesisGenitive{\ifEN bachelor \else bakalářské \fi}
%\def\ThesisGenitive{\ifEN master \else diplomové \fi}
%\def\ThesisGenitive{\ifEN rigorous \else rigorózní \fi}
%\def\ThesisGenitive{\ifEN doctoral \else disertační \fi}
\def\ThesisAccusative{\ifEN bachelor \else bakalářskou \fi}
%\def\ThesisAccusative{\ifEN master \else diplomovou \fi}
%\def\ThesisAccusative{\ifEN rigorous \else rigorózní \fi}
%\def\ThesisAccusative{\ifEN doctoral \else disertační \fi}



%%% Fill in your details %%%

% (Note: \xxx is a "ToDo label" which makes the unfilled visible. Remove it.)
\def\ThesisTitle{Možnosti návrhového vzoru Entity-Component-System: případová studie}
\def\ThesisAuthor{Petr Kotáb}
\def\YearSubmitted{2024}

% department assigned to the thesis
\def\Department{Katedra softwaru a výuky informatiky}
% Is it a department (katedra), or an institute (ústav)?
\def\DeptType{Katedra}

\def\Supervisor{Mgr. Pavel Ježek, Ph.D.}
\def\SupervisorsDepartment{Katedra distribuovaných a spolehlivých systémů}

% Study programme and specialization
\def\StudyProgramme{Informatika }
\def\StudyBranch{Počítačová grafika, vidění a vývoj her}

\def\Dedication{%
Dedication. \xxx{It is nice to say thanks to supervisors, friends, family, book authors and food providers.}
}

\def\AbstractEN{%
\xxx{Abstracts are an abstract form of art. Use the most precise, shortest sentences that state what problem the thesis addresses, how it is approached, pinpoint the exact result achieved, and describe the applications and significance of the results. Highlight anything novel that was discovered or improved by the thesis. Maximum length is 200 words, but try to fit into 120. Abstracts are often used for deciding if a reviewer will be suitable for the thesis; a well-written abstract thus increases the probability of getting a reviewer who will like the thesis.}
% ABSTRACT IS NOT A COPY OF YOUR THESIS ASSIGNMENT!
}

\def\AbstractCS{%
\xxx{You will need to submit both Czech and English abstract to the SIS, no matter what language you use in the thesis. If writing in English, translate the contents of \texttt{\textbackslash{}AbstractEN} into this field. In case you do not speak czech, your supervisor should be able to help you with the translation.}
}

% 3 to 5 keywords (recommended), each enclosed in curly braces.
% Keywords are useful for indexing and searching for the theses by topic.
\def\Keywords{
    {návrhový vzor}
    {Entity-Component-System}
    {počítačové hry}
    {simulace}
    {případová studie}
}

% If your abstracts are long and do not fit in the infopage, you can make the
% fonts a bit smaller by this setting. (Also, you should try to compress your abstract more.)
% Alternatively, consider increasing the size of the page by uncommenting the
% geometry modification in thesis.tex.
\def\InfoPageFont{}
%\def\InfoPageFont{\small}  %uncomment to decrease font size

\ifEN\relax\else
% If you are writing a czech thesis, you additionally need to fill in the
% english translation of the metadata here!
\def\ThesisTitleEN{Exploring Options of Entity-Component-System Design Pattern: A Case Study}
\def\DepartmentEN{Department of software and computer science education}
\def\DeptTypeEN{Department}
\def\SupervisorsDepartmentEN{Department of Distributed and Dependable Systems}
\def\StudyProgrammeEN{Computer Science}
\def\StudyBranchEN{Computer Graphics, Vision and Game Development}
\def\KeywordsEN{
    {design pattern}
    {Entity-Component-System}
    {computer games}
    {simulation}
    {case study}
}
\fi


\usepackage[a-2u]{pdfx}

\ifEN\else\usepackage[czech,shorthands=off]{babel}\fi

% See https://en.wikipedia.org/wiki/Canons_of_page_construction before
% modifying the size of printable area. LaTeX defaults are great.
% If you feel it would help anything, you can enlarge the printable area a bit:
%\usepackage[textwidth=390pt,textheight=630pt]{geometry}
% The official recommendation expands the area quite a bit (looks pretty harsh):
%\usepackage[textwidth=145mm,textheight=247mm]{geometry}

%%% TYPICAL FONT CHOICES (uncomment what you like) %%%
% Recommended combo: Libertinus (autoselects Biolinum for sans) and everything
% else (math+tt) comes from Latin Modern)
\usepackage{lmodern}
\usepackage[mono=false]{libertinus}

% For the "classic" LaTeX fonts (very good for pure math theses), simply
% comment out the libertinus package above.

% IBM Plex font suite: nice, but requires us to fine-tune the sizes and does
% not directly support small caps (\textsc):
%\usepackage[usefilenames,RM={Scale=0.88},SS={Scale=0.88},SScon={Scale=0.88},TT={Scale=0.88},DefaultFeatures={Ligatures=Common}]{plex-otf}

% TeX Gyre combo (Pagella+Heros+Cursor)
%\usepackage{fontspec}
%\setmainfont{TeX Gyre Pagella}
%\setsansfont{TeX Gyre Heros}
%\setmonofont{TeX Gyre Cursor}

% some useful packages
\usepackage{microtype}
\usepackage{amsmath,amsfonts,amsthm,bm}
\usepackage{graphicx}
\usepackage{xcolor}
\usepackage{booktabs}
\usepackage{caption}
\usepackage{floatrow}

% Bibliography formatting.
% CHECK THE REQUIREMENTS OF YOUR DEPARTMENT AND FACULTY ON THE CITATION FORMAT!
%
% These are relatively "safe" default options that most people use:
\usepackage[natbib,style=numeric,sorting=none]{biblatex}
% alternative with alphanumeric citations (more informative than numbers, and
% more common in computer science journals):
%\usepackage[natbib,style=alphabetic]{biblatex}
%
% ALTERNATIVES THAT CONFORM TO ISO690
% ISO690 is not the greatest citation format ever, but may be formally
% required at Charles University, depending on your faculty and department.
%\usepackage[natbib,style=iso-numeric,sorting=none]{biblatex}
%\usepackage[natbib,style=iso-alphabetic]{biblatex}
% You might want to add extra options such as `maxbibnames=6,maxcitenames=2`
% here to further conform to some of the formatting requirements (see below for
% details). Again, consult your faculty rules.
%
% Additional option choices:
%  - add `giveninits=true` to typeset "E. A. Poe" instead of full Edgar Allan
%  - `terseinits=true` additionaly shortens it to nature-like "Poe EA"
%  - add `maxnames=10` to limit (or loosen) the maximum number of authors in
%    bibliography entry before shortening to `et al.` (useful when referring to
%    book collections that may have hundreds of authors)
%  - use `maxcitenames=2` to finetune the amount of authors listed in text-cite
%    commands (\citet). Corresponding option that only affects the bibliography
%    is `maxbibnames=10`.
%  - `sorting=none` causes the bibliography list to be ordered by the order of
%    citation as they appear in the text, which is usually the desired behavior
%    with numeric citations. Additionally you can use a style like
%    `numeric-comp` that compresses the long lists of citations such as
%    [1,2,3,4,5,6,7,8] to simpler [1--8]. This is especially useful if you plan
%    to add tremendous amounts of citations, as usual in life sciences and
%    bioinformatics.
%  - if you don't like the "In:" appearing in the bibliography, use the
%    extended style (`ext-numeric` or `ext-alphabetic`), and add option
%    `articlein=false`.
%
% possibly reverse the names of the authors with the default styles:
%\DeclareNameAlias{default}{family-given}

% load the file with bibliography entries
\addbibresource{refs.bib}

% remove this if you won't use fancy verbatim environments
\usepackage{fancyvrb}

% remove this if you won't typeset TikZ graphics
\usepackage{tikz}
\usetikzlibrary{positioning} %add libraries as needed (shapes, decorations, ...)

% remove this if you won't typeset any pseudocode
\usepackage{algpseudocode}
\usepackage{algorithm}

% remove this if you won't list any source code
\usepackage{listings}


\hypersetup{unicode}
\hypersetup{breaklinks=true}

\usepackage[noabbrev]{cleveref}


% various forms of TODOs (you should remove this before submitting)
\usepackage[textsize=tiny, backgroundcolor=yellow!25, linecolor=black!25]{todonotes}
\newcommand{\xxx}[1]{\textcolor{red!}{#1}}

 % remove this before compiling the final version


% use this for typesetting a chapter without a number, e.g. intro and outro
\def\chapwithtoc#1{\chapter*{#1}\addcontentsline{toc}{chapter}{#1}}

% If there is a line/figure overflowing into page margin, this will make the
% problem evident by drawing a thick black line at the overflowing spot. You
% should not disable this.
\overfullrule=3mm

% The maximum stretching of a space. Increasing this makes the text a bit more
% sloppy, but may prevent the overflows by moving words to next line.
\emergencystretch=1em

\ifEN
\theoremstyle{plain}
\newtheorem{thm}{Theorem}
\newtheorem{lemma}[thm]{Lemma}
\newtheorem{claim}[thm]{Claim}
\newtheorem{defn}{Definition}
\theoremstyle{remark}
\newtheorem*{cor}{Corollary}
\else
\theoremstyle{plain}
\newtheorem{thm}{Věta}
\newtheorem{lemma}{Lemma}
\newtheorem{claim}{Tvrzení}
\newtheorem{defn}{Definice}
\theoremstyle{remark}
\newtheorem*{cor}{Důsledek}
\fi

\newenvironment{myproof}{
  \par\medskip\noindent
  \textit{\ifEN Proof \else Důkaz \fi}.
}{
\newline
\rightline{$\qedsymbol$}
}

% real/natural numbers
\newcommand{\R}{\mathbb{R}}
\newcommand{\N}{\mathbb{N}}

% asymptotic complexity
\newcommand{\asy}[1]{\mathcal{O}(#1)}

% listings and default lstlisting config (remove if unused)
\DeclareNewFloatType{listing}{}
\floatsetup[listing]{style=ruled}

\DeclareCaptionStyle{thesis}{style=base,font={small,sf},labelfont=bf,labelsep=quad}
\captionsetup{style=thesis}
\captionsetup[algorithm]{style=thesis,singlelinecheck=off}
\captionsetup[listing]{style=thesis,singlelinecheck=off}

% Customization of algorithmic environment (comment style)
\renewcommand{\algorithmiccomment}[1]{\textcolor{black!25}{\dotfill\sffamily\itshape#1}}

% Uncomment for table captions on top. This is sometimes recommended by the
% style guide, and even required for some publication types.
%\floatsetup[table]{capposition=top}
%
% (Opinionated rant:) Captions on top are not "compatible" with the general
% guideline that the tables should be formatted to be quickly visually
% comprehensible and *beautiful* in general (like figures), and that the table
% "head" row (with column names) should alone communicate most of the content
% and interpretation of the table. If you just need to show a long boring list
% of numbers (because you have to), either put some effort into showing the
% data in an attractive figure-table, or move the data to an attachment and
% refer to it, so that the boredom does not impact the main text flow.
%
% You can make the top-captions look much less ugly by aligning the widths of
% the caption and the table, with setting `framefit=yes`, as shown below.  This
% additionally requires some extra markup in your {table} environments; see the
% comments in the example table in `ch2.tex` for details.
%\floatsetup[table]{capposition=top,framefit=yes}

\ifEN\floatname{listing}{Listing}
\else\floatname{listing}{Výpis kódu}\fi
\lstset{ % use this to define styling for any other language
  language=C++,
  tabsize=2,
  showstringspaces=false,
  basicstyle=\footnotesize\tt\color{black!75},
  identifierstyle=\bfseries\color{black},
  commentstyle=\color{green!50!black},
  stringstyle=\color{red!50!black},
  keywordstyle=\color{blue!75!black}}

% Czech versions of the used cleveref references (It's not as convenient as in
% English because of declension, cleveref is limited to sg/pl nominative. Use
% plain \ref to dodge that.)
\ifEN\relax\else
\crefname{chapter}{kapitola}{kapitoly}
\Crefname{chapter}{Kapitola}{Kapitoly}
\crefname{section}{sekce}{sekce}
\Crefname{section}{Sekce}{Sekce}
\crefname{subsection}{sekce}{sekce}
\Crefname{subsection}{Sekce}{Sekce}
\crefname{subsubsection}{sekce}{sekce}
\Crefname{subsubsection}{Sekce}{Sekce}
\crefname{figure}{obrázek}{obrázky}
\Crefname{figure}{Obrázek}{Obrázky}
\crefname{table}{tabulka}{tabulky}
\Crefname{table}{Tabulka}{Tabulky}
\crefname{listing}{výpis}{výpisy}
\Crefname{listing}{Výpis}{Výpisy}
\floatname{algorithm}{Algoritmus}
\crefname{algorithm}{algoritmus}{algoritmy}
\Crefname{algorithm}{Algoritmus}{Algoritmy}
\newcommand{\crefpairconjunction}{ a~}
\newcommand{\crefrangeconjunction}{ a~}
\fi
 % use this file for various custom definitions

\begin{document}

% the layout is mandatory, edit only in dire circumstances

\pagestyle{empty}
\hypersetup{pageanchor=false}
\begin{center}

% top part of the layout, this actually differs between faculties

\def\ThesisTitleXmff{%
  \ifEN
    \centerline{\mbox{\includegraphics[width=166mm]{img/logo-en.pdf}}}
  \else
    \centerline{\mbox{\includegraphics[width=166mm]{img/logo-cs.pdf}}}
  \fi
  \vspace{-8mm}\vfill%
  {\bf\Large\ThesisTypeTitle}
  \vfill%
  {\LARGE\ThesisAuthor}\par
  \vspace{15mm}%
  {\LARGE\bfseries\ThesisTitle}
  \vfill%
  \Department}
\def\ThesisTitleCuniLogo#1{%
  {\large\UKName\par\medskip\par\UKFaculty }
  \vfill%
  {\bf\Large\ThesisTypeTitle}
  \vfill%
  \includegraphics[width=70mm]{#1}
  \vfill%
  {\LARGE\ThesisAuthor}\par
  \vspace{15mm}%
  {\LARGE\bfseries\ThesisTitle}
  \vfill%
  \Department\par}
\def\ThesisTitleXcuni{\ThesisTitleCuniLogo{img/uklogo.pdf}}
\def\ThesisTitleXnatur{\ThesisTitleCuniLogo{img/naturlogo.pdf}}

% choose the correct page and print it
\csname ThesisTitleX\ThesisTitleStyle\endcsname
% latex corner: X is the new @

\vfill

{
\centerline{\vbox{\halign{\hbox to 0.45\hsize{\hfil #}&\hskip 0.5em\parbox[t]{0.45\hsize}{\raggedright #}\cr
\ifEN Supervisor of the \ThesisGenitive thesis:
\else Vedoucí \ThesisGenitive práce: \fi
& \Supervisor \cr
\noalign{\vspace{2mm}}
\ifEN Study programme: \else Studijní program: \fi
& \StudyProgramme \cr
\noalign{\vspace{2mm}}
\ifEN Study branch: \else Studijní obor: \fi
& \StudyBranch \cr
}}}}

\vfill

\ifEN Prague \else Praha \fi
\YearSubmitted

\end{center}

\newpage

% remember to sign this!
\openright
\hypersetup{pageanchor=true}
\pagestyle{plain}
\pagenumbering{roman}
\vglue 0pt plus 1fill

\ifEN
\noindent
I declare that I carried out this \ThesisAccusative thesis independently, and only with the cited
sources, literature and other professional sources. It has not been used to obtain another
or the same degree.
\else
\noindent
Prohlašuji, že jsem tuto \ThesisAccusative práci vypracoval(a) samostatně a výhradně
s~použitím citovaných pramenů, literatury a dalších odborných zdrojů.
Tato práce nebyla využita k získání jiného nebo stejného titulu.
\fi

\ifEN
\medskip\noindent
I understand that my work relates to the rights and obligations under the Act No.~121/2000 Sb.,
the Copyright Act, as amended, in particular the fact that the Charles
University has the right to conclude a license agreement on the use of this
work as a school work pursuant to Section 60 subsection 1 of the Copyright~Act.
\else
\medskip\noindent
Beru na~vědomí, že se na moji práci vztahují práva a povinnosti vyplývající
ze zákona č. 121/2000 Sb., autorského zákona v~platném znění, zejména skutečnost,
že Univerzita Karlova má právo na~uzavření licenční smlouvy o~užití této
práce jako školního díla podle §60 odst. 1 autorského zákona.
\fi

\vspace{10mm}


\ifEN
\hbox{\hbox to 0.5\hsize{%
In \hbox to 6em{\dotfill} date \hbox to 6em{\dotfill}
\hss}\hbox to 0.5\hsize{\dotfill\quad}}
\smallskip
\hbox{\hbox to 0.5\hsize{}\hbox to 0.5\hsize{\hfil Author's signature\hfil}}
\else
\hbox{\hbox to 0.5\hsize{%
V \hbox to 6em{\dotfill} dne \hbox to 6em{\dotfill}
\hss}\hbox to 0.5\hsize{\dotfill\quad}}
\smallskip
\hbox{\hbox to 0.5\hsize{}\hbox to 0.5\hsize{\hfil Podpis autora\hfil}}
\fi

\vspace{20mm}
\newpage

% dedication

\openright

\noindent
\Dedication

\newpage

% mandatory information page

\openright

\vbox to 0.49\vsize{\InfoPageFont
\setlength\parindent{0mm}
\setlength\parskip{5mm}

\ifEN Title: \else Název práce: \fi
\ThesisTitle

\ifEN Author: \else Autor: \fi
\ThesisAuthor

\DeptType:
\Department

\ifEN Supervisor: \else Vedoucí \ThesisGenitive práce: \fi
\Supervisor, \SupervisorsDepartment

\ifEN Abstract: \AbstractEN \else Abstrakt: \AbstractCS \fi

\ifEN Keywords: \else Klíčová slova: \fi
\Keywords

\vss}\ifEN\relax\else\nobreak\vbox to 0.49\vsize{\InfoPageFont
\setlength\parindent{0mm}
\setlength\parskip{5mm}

Title:
\ThesisTitleEN

Author:
\ThesisAuthor

\DeptTypeEN:
\DepartmentEN

Supervisor:
\Supervisor, \SupervisorsDepartmentEN

Abstract:
\AbstractEN

Keywords:
\KeywordsEN

\vss}
\fi

\newpage

\openright
\pagestyle{plain}
\pagenumbering{arabic}
\setcounter{page}{1}


\tableofcontents

%\setlength{\parskip}{\baselineskip}

\chapter{Úvod}

Vytváření her je komplexní proces, který zahrnuje mnoho fází, včetně prototypování a iterací. Často se ocitneme v situaci, kdy je nezbytné výrazně upravit již existující kód. Z tohoto důvodu je klíčová vysoká flexibilita architektury kódu.

Jedním z přístupů k řešení tohoto problému je využití návrhového vzoru Component, který poskytuje vysokou míru flexibility. Tento návrhový vzor je běžně používán v herních frameworcích a enginech, jako je například Unity (koncept MonoBehaviour) nebo MonoGame (koncept GameComponent). Návrhový vzor Component je dobře představen v knize Game Programming Patterns~\cite{nystrom2014game}.

Nicméně existují pokročilejší přístupy, které mohou nabídnout ještě lepší flexibilitu a výkon. Jedním takovím je návrhový vzor Entity-Component-System (ECS).  I přesto, že se v současné době nepoužíva tak často jako návrhový vzor Component, tak nabízí značné výhody při tvorbě her. Proto se na něj v této práci zaměříme. Tento návrhový vzor je dobře rozebrán v sérii článků s názvem ECS back and forth~\cite{Caini_2019}.

\section{Entity-Component-System}
ECS je návrhový vzor používaný pro vývoj her. Jeho základními stavebními kameny jsou entity, komponenty a systémy.

\begin{enumerate}
    \item \textbf{Entita (Entity):} Reprezentuje objekt v herní scéně, může být cokoliv od hráče po strom nebo zvíře.
    \item \textbf{Komponenta (Component):} Entita je složena z komponent, které ji popisují. Komponenty obsahují pouze data a žádnou herní logiku. Například hráč může mít komponentu Control, která popisuje, jak se ovládá, strom může mít komponentu Resource, která popisuje, kolik dřeva hráč získá, pokud jej pokácí, a zvíře může mít komponentu PathFollow, která popisuje cestu, po které se pohybuje. Všechny tři entity mohou mít komponenty Appearance, která popisuje, jak daný objekt vypadá, a Position, popisující, kde se daný objekt nachází v herní scéně. Dále hráč a zvíře mohou mít komponentu Movement, popisující jejich pohyb.
    \item \textbf{Systém (System):} Obsahuje herní logiku. Systémy pracují nad určitou množinou komponent. V každé iteraci vezmou všechny entity obsahující tuto množinu komponent a provedou nad nimi určité operace. Například:
    \begin{itemize}
        \item \textbf{InputSystem:} Pracuje nad entitami s komponentami Control a Position, řídí logiku jejich ovládání.
        \item \textbf{PathMovementSystem:} Pracuje nad entitami s komponentami PathFollow a Movement, řídí logiku jejich pohybu po cestě.
        \item \textbf{RenderSystem:} Pracuje nad entitami s komponentou Appearance a řídí logiku jejich vykreslení.
    \end{itemize}
\end{enumerate}

ECS nejenže poskytuje vysokou flexibilitu, ale také vyniká v oblasti výkonu. Jeho flexibilitu a výkon si více rozebereme v~\cref{chap:ecs} \xxx{poresit sklonovani u linku}. Mezi cíle této práce bude patřit prozkoumat možnosti, které tento návrhový vzor přináší, a ukázat výhody, které nám ECS může poskytnout v porovnání s návrhovým vzorem Component.

\section{ECS knihovny}
Aby bylo možné efektivně využívat návrhový vzor ECS, je nezbytné implementovat odpovídající infrastrukturu. Samotná implementace této infrastruktury představuje obtížnou a časově náročnou úlohu, což vede mnoho vývojářů k volbě již existujících knihoven implementujících infrastrukturu pro tento návrhový vzor.

V této práci se zaměříme na knihovny pro programovací jazyk C\#, neboť se jedná o programovací jazyk běžně používaný pro hry a aplikace. I když se omezíme pouze na C\#, existuje velké množství těchto knihoven.

Jedním ze způsobů jak si vývojáři mohou vybrat ECS knihovnu je použít porovnání. CSharpECSComparison~\cite{CSharpECSComparison} je projekt na platformě GitHub, který porovnává populární ECS knihovny pro C\# na základě funkcionalit, vlastností a závislostí. Vývojář by si tedy mohl říci, které vlastnosti a funkcionality by od takové knihovny požadoval a na základě toho si vybrat. Ovšem možnosti co nám jednotlivé ECS knihovny nabízejí se tolik neliší, proto dává smysl porovnávat jednotlivé knihovny na základě výkonu. Bohužel CSharpECSComparison nám neposkytuje žádné informace o výkonu jednotlivých knihoven.

Ecs.CSharp.Benchmark~\cite{EcsCsharpBenchmark} je projekt na platforme GitHub, který provádí jednoduché výkonnostní testy ECS knihoven pro C\#. Výsledky těchto testů mohou být užitečné při výběru knihovny na základě výkonu. Problém však může spočívat v tom, že tyto testy jsou příliš jednoduché. Každý test, který Ecs.CSharp.Benchmark provádí, pracuje pouze s velmi jednoduchými komponentami, kde každá komponenta obsahuje pouze jedno celé číslo. Testy jsou rozděleny do několika sad. První sada vytváří entity s těmito komponentami. Druhá sada spouští jednoduché systémy, které pouze sčítají výše zmíněná celá čísla na těchto komponentách.

Jelikož jsou testy příliš jednoduché, v případě rozsáhlejšího projektu, jako je hra, by se výsledky mohly lišit. Jednotlivé entity mohou mít velké množství komponent, z nichž každá může být složitější. Systémy mohou být také mnohem komplexnější a provádět mnohem složitější úkoly než pouhé sčítání celých čísel.

Jedním ze způsobů, jak by bylo možné lépe porovnat výkon jednotlivých ECS knihoven, by bylo vytvoření hry, na které by bylo možné měřit výkon jednotlivých knihoven. Taková hra by musela být nezávislá na konkrétní ECS knihovně. Měla by také obsahovat netriviálně velký svět s větším počtem entit. To by nám umožnilo lépe provádět měření. Jedním z cílů této práce bude vytvoření takovéto hry.

Poté, co budeme mít tuto hru k dispozici, bude možné vytvořit výkonnostní test, který by dokázal změřit výkon příslušné ECS knihovny na této hře. Pomocí tohoto testu bude možné udělat výkonostní porovnání jednotlivých ECS knihoven. Vytvoření tohoto porovnání bude dalším cílem této práce.

Zároveň takto vytvořená hra pro nás bude velice užitečná neboť nám pomůže při průzkumu možností návrhové vzoru ECS, což je jedním z cílů této práce.

\section{Abstrakční vrstva}
Vzhledem k tomu, že plánovanou hru budeme chtít využít k porovnávání různých ECS knihoven, je nezbytné navrhnout ji tak, aby fungovala nezávisle na konkrétní ECS knihovně. S ohledem na odlišné API jednotlivých ECS knihoven je nezbytné vytvořit abstrakční vrstvu pro ECS. Pro každou knihovnu, kterou budeme chtít porovnávat, poskytneme implementaci této vrstvy. Samotná hra poté pracuje s touto abstrakční vrstvou.

Vytvoření takovéto vrstvy představuje netriviální úkol, navzdory podobnosti rozhraní jednotlivých ECS knihoven. Významný rozdíl spočívá v rozhraní systémů, které se výrazně liší mezi jednotlivými knihovnami.

\section{Cíle práce}
Pro shrnutí, tato práce bude mít následující cíle:
\begin{enumerate}
    \item[h1)] Vytvořit hru pomocí ECS, která bude nezávislá na konkrétní ECS knihovně.
    \item[h2)] Pomocí vytvořené hry porovnat relativní výkon populárních ECS knihoven pro programovací jazyk C\#.
\end{enumerate}
S využitím této hry také budeme chtít:
\begin{enumerate}
    \item[v1)] Prozkoumat možnosti návrhového vzoru ECS.
    \item[v2)] Ukázat výhody, které nám ECS může nabídnout oproti návrhovému vzoru Component.
\end{enumerate}














\iffalse
\chapwithtoc{Introduction}

Introduction should answer the following questions, ideally in this order:
\begin{enumerate}
\item What is the nature of the problem the thesis is addressing?
\item What is the common approach for solving that problem now?
\item How this thesis approaches the problem?
\item What are the results? Did something improve?
\item What can the reader expect in the individual chapters of the thesis?
\end{enumerate}

Expected length of the introduction is between 1--4 pages. Longer introductions may require sub-sectioning with appropriate headings --- use \texttt{\textbackslash{}section*} to avoid numbering (with section names like `Motivation' and `Related work'), but try to avoid lengthy discussion of anything specific. Any ``real science'' (definitions, theorems, methods, data) should go into other chapters.
\todo{You may notice that this paragraph briefly shows different ``types'' of `quotes' in TeX, and the usage difference between a hyphen (-), en-dash (--) and em-dash (---).}

It is very advisable to skim through a book about scientific English writing before starting the thesis. I can recommend `\citetitle{glasman2010science}' by \citet{glasman2010science}.
\fi
\chapter{Architektury her}
V této kapitole si rozebereme architektury, které lze využít pro vývoj her. Pomůže nám to lépe porozumět problémům, které ECS řeší. Začneme objektovým návrhem a prozkoumáme problémy spojené s dědičností při vývoji her. Následně se zaměříme na návrhový vzor Component, který se snaží řešit tuto problematiku, a nahlédneme proč komunikace mezi komponentami může být problematická.

% Iterativni proces
Proto abychom lépe porozuměli herním architekturám, bude nutné si nejprve přiblížit jak se tvorba her liší od tvorby jiného typu softwaru. Vývoj her je iterativní proces, při kterém se nejprve navrhnou herní mechaniky. Následně probíhá testovaní, při kterém se zjistí, jestli jsou příslušné herní mechaniky zábavné. Pokud nejsou je nutné je upravit a znovu otestovat. Tyto dvě aktivity se opakují dokud vytvořené mechaniky nejsou dostatečně zábavné. V moment co jsou mechaniky vytvořené a víme o nich, že jsou zábavné, je možné vytvořit další mechaniky a celí proces opakovat. Takto se postupuje dokud nedojde k vytvoření celé hry.

% Co to implikuje?
Výše zmíněný iterativní proces, nám implikuje to, že zdrojový kód pro hry se mění velice často. Z toho důvodu je flexibilita velice důležitá vlastnost architektury při vývoji her.

\section{Objektový návrh}
% prirozenym zpusobem vyuzit bezneho objektoveho navrhu programovaciho jazyka, cili udelat nejakou sadu objektu provazanych dedicnosti
Přirozeným přístupem pro tvorbu software je použítí objektový návrh a s jeho pomocí si vytvořit sadu tříd provázaných dědičností. Ovšem jak již bylo naznačeno, v případě tvorby her může být tento přístup problematický.

% Jak se pouziva objektovy nabrh ve vyvoji her

Navrhneme si jednoduchou hru, na které si přiblížíme problémy s dědičností, které mohou při vývoji her s objektovým návrhem nastat. Budeme navrhovat jednoduchý tower defense \xxx{link}. Začneme třídou \verb|GameObject|, ta nám bude sloužit jako předek pro libovolný objekt ve scéně. Nyní bychom chtěli vytvořit jednotky, vytvoříme tedy třídu \verb|Unit|, která podědí od \verb|GameObject|. Chtěli bychom aby ve hře byli také speciální jednotky, které umí střílet, např. lučišníci. Vytvoříme tedy třídu \verb|Archer|, která dědí od třídy \verb|Unit|. Dále v naší hře budem potřebovat věže, uděláme tedy třídu \verb|Tower|, dědicí od třídy \verb|GameObject|. Určitě budem chtít od většiny našich věží aby uměli střílet, implementujeme tedy další třídu \verb|ShootingAbleTower|, dědicí od třídy \verb|Tower|. Do této třídy budem určitě chtít implementovat mechaniku střelby. Ovšem střelbu jsme již určitě implementovali do třídy \verb|Archer|, což nám nyní vytváří duplicitní kód.

\xxx{obrazek}

Jak tento problém vyřešit? Prvním přístupem je využití multiple inheritance \xxx{link} (vícenásobné dědičnosti). Pokud programovací jazyk povoluje multiple inheritance, umožňuje to jedné třídě mít více než jednoho předka. Pokud bychom chtěli pomocí multiple inheritence vyřešit náš problém, mohli bychom implementovat třídu \verb|ShootingAble|, která dědí od \verb|GameObject|. Do této třídy bychom implementovali mechaniku střelby. Náš \verb|Archer| by poté dědil od \verb|ShootingAble| a od \verb|Unit|. Naše \verb|ShootingAbleTower| by dědila od \verb|ShootingAble| a od \verb|Tower|.

\xxx{obrazek}

Tímto řešení by nám ale vznikl tzv. diamond problém \xxx{link}. Jedná se problém, který může vzniknout při používání multiple inheritance. Diamond problém nástává tehdy, pokud je potřeba přisoupit ke členům definovaných na společném předkovi. V našem případě se jedná o členy na třídě \verb|GameObject|.

Alternativním přístupem by mohlo být přesunutí společné logiky do společného předka. V našem případě bychom mechaniku střelby implementovali přímo ve třídě \verb|GameObject|.

Ovšem tímto porušíme jeden z principů objektového programování, konkrétně single responsibility \xxx{link}. Tento princip nám říká, že každá třída by měla být zodpovědná pouze za jednu věc. Porušení tohoto principu může mít za následek, že se nám ve společném předkovi začne hromadit velké množství logiky. To může vyústit v obrovskou třídu, která řeší všechny možné herné mechaniky. Jedná se o tzv. blob object \xxx{link}. 

Další problém je, že náš návrh je příliš závislý na hierarchii dědičnosti, což omezuje design naší hry. V praxi to znamená, že např. není možné za běhu přidat objektu schopnost střílet.

Řešením celé této situace je koncept composition over inheritance \xxx{link}. Tento princip nám říka, že namísto hluboké hierarchie dědičnosti bychom měli preferovat rozdělení objektů do komponent. Je to princip, na kterém je postavený návrhový vzor Component a také ECS.

\xxx{Pridat jako poznamku pod carou: Ve skutečnosti byl princip composition over inheritance zmíněný již v Gang of Four a jedná se o dobrou praktiku v objektovém programování. }

\section{Návrhový vzor Component}
\label{sec:component}
Component je návrhový vzor... .

Tento návrhový vzor je velice populární.

Nyní si rozeberme problemy, které mohou vzniknout při používání návrhového vzoru Component.

Příklad?
\chapter{ECS}
\label{chap:ecs}
V této kapitole si více rozeberem návrhový vzor ECS. Řekneme si jak se implementuje ukládání komponent, rozebereme si vlastnosti ECS a na závěr porovnáme ECS s návrhovým vzorem Component.

\section{Reprezentace ECS}
To jak ECS funguje jsme si již rozebrali v úvodu. Nyní si rozebereme jak bývají jednotlivé prvky ECS reprezentovány při jejich implementaci. Krom základních členů, jako jsou entity, komponenty a systémy, si představíme ještě world a query.

Komponenty, jako takové, bývají často reprezentovány třídou nebo strukturou. Jak již bylo zmíněno v úvodu, komponenty obsahují pouze data, nikoliv žádnou logiku. Proto tyto třídy a struktury neobsahují žádné funkce.

Entita by, po vzoru návrhového vzoru Component, mohla být reprezentována jako kolekce svých komponent. Ovšem to se v praxi nedělá, namísto toho se používá řešení, které jak si později ukážeme vede k lepšímu výkonu. Běžné řešení bývá to, že každá ECS implementace má správce všech entit a komponent, ve kterém jsou jednotlivé komponenty uloženy. Jednotlivé entity jsou potom reprezentovány pouze jako jednoduchý identifikátor a je úlohou tohoto správce aby mapoval jednotlivé entity k jejich komponentám.

Tomuto správce se v ECS implementacích běžně říká world. Jak již bylo zmíněno, slouží jako správce a kontejner pro všechny entity a komponenty v herním světě. Jeho rozhraní nabízí funkce, pomocí niž je možné vytvářet a mazat jednotlivé entity. Často obsahuje také funkce pro přidávání a odebírání komponent jednotlivím entitám.

Některé ECS implemetace reprezentují entity pomocí malé struktury, obsahujicí identifikátor dané entity společně s referencí na world do kterého patří. To umožnuje mít této struktuře bohaté rozhraní s funkcemi pro přidávání a odebírání komponent.

Reprezentace systémů se v jednotlivích ECS implementacích dost liší. Některé vyžadují aby systém byl třída dědicí od abstraktního předka, jíné zase volí opačný extrém a dovolují systém reprezentovat téměř jakkoliv. Navzdory velkým odlišnostem se ve velkém počtu ECS implementací často objevuje objekt, který mohou používat jednotlivé systémy pro iterování entity s určitou množinou komponent.

Tomuto objektu se v ECS implementacích běžně říká query. Pro jeho vytvoření je většinou nutné poskytnou informaci, nad jakými typy komponent bude query pracovat. Je například možné mít Query, které pracuje nad komponenty Movement a Position. Takové Query by pak bylo schopné iterovat nad všemi entitami co mají komponenty Movement a Position. Pomocí tohoto Query by bylo možné implementovat MovementSystem, který by řídil logiku pohybu.

Query se dělí do dvou typů, statické a dynamické. Statické query je možné zkonstruovat (znamenaje specifikovat typy komponent nad kterými bude pracovat) pouze za kompilace. Naproti tomu dynamické query je možné zkonstruovat za běhu. Dynamické query je poměrně pokročilá věc a pouze malé množství ECS implementací jej podporuje.

Je důležité si uvědomit, že správa entit a komponent je dost složitá úloha. Je nutné aby bylo možné namapovat jednotlivé komponenty na jejich entity. Zároveň je také zapotřebí aby bylo možné iterovat určité n-tice těchto komponent za pomocí query. Je také důležité brát v potaz to, že jednotlivé entity a komponenty se mohou přidávat a odebírat za běhu. Více si o těchto věcech řekneme v pozdější sekci, nyní se ale zamřme na mazání entit. Jelikož smazání takové entity může být naročné na implementaci a nebo také na výkon, některé ECS implementace využívajicí koncept tombstone. Tombstone je oznaceni pro objekt, ktery po smazani stale existuje, pouze je někde poznamenáno, že je již smazaný. Smazané entity bývají tedy v některých ECS implementacích označeny jako tombstone a query je při iteraci ignoruje.

%\section{Správa paměti}
%Než si rozeberem jak je možné úkládat komponenty, bude nutné si přiblížit jak vlastně funguje cache v počítači a také jaký je rozdíl mezi hodnotovími a referenčními typy v C\#.

%Cache

%Referencni vs hodnotove typy

\xxx{sekce o cache a referencni vs hodnotove typy v C\#?}

\section{Ukladaní komponent}
Nyní si rozeberem způsoby, kterými jednotlivé ECS implementace řeší ukládání komponent v paměti. Používají se převážně dva způsoby arch type a sparse set. Ovšem některé knihovny tyto způsoby nepoužívají a ukládají komponenty neefektivně, přiblížíme si tedy také jak to tyto knihovny dělají a proč je to neefektivní.

Arch type

Sparse set

Componenty jako pointery

\section{Vlastnosti ECS}
...

Vysoka flexibilita

Vysoky vykon

Flexibilita > vykon

\section{ECS vs Component}

\xxx{TODO}

\xxx{maybe topics: relace, eventy, tag component, system types, reaction systems, paralelism, paralel queries, query filters, functional programming, relace, eventy, prukopnici (herni engine Bevy, ECS knihovny entitas a flecs)}
\chapter{Co požadujeme od hry}
V této kapitole si rozebereme požadavky, které požadujeme od naší hry. Prvně si rozebereme jednotlivé charakteristiky, které by naše hra měla splňovat. Dále na základě těchto charakteristik zvolíme žánr naší hry a odůvodníme proč jsme jej zvolili.

Než budeme pokračovat, připomeňme si prvně, že naší hru chceme použít pro měření výkonu ECS knihoven. Budeme tedy chtít aby byla nezávislá na konkrétní knihovně. Konkrétní knihovnu bude možné zvolit před spuštěním hry.

\section{Charakteristiky hry}
Aby nám naše hra napomohla kde splnění práce, budeme chtít aby splňovala následující charakteristiky:

\begin{enumerate}
    \item Nezávislost na konkrétní ECS knihovně.
    \item Hra by neměla využívat častého přidávání a odebírání komponent.
    \item Velký počet entit.
    \item Adekvátní složitost hry.
\end{enumerate}

Nezávislost na konkrétní ECS knihovně pro nás bude důležitá, abychom hru mohli použít pro měření výkonu jednotlivých knihoven. Hra bude namísto konkrétní ECS knihovny pracovat pouze s abstrakční vrstvou. Bude pro nás žádoucí aby se výkon naší hry bežící s konkretní ECS knihovnou, implementující abstrakční vrstvu, blížil k výkonu hry, která by používala pouze konkrétní ECS knihovnu bez abstrakční vrstvy.

Časté přidávání a odebírání komponent by vedlo ke zvýhodnění ECS knihoven založených na sparse set. Cílem této práce ale není porovnat knihovny založené na sparse set oproti knihovnám založených na arch type. To jestli hra využívá častého přidávání a odebírání komponent je závislé čistě na implementaci. 

Ukažme si to na příkladu, konkrétně bychom chtěli reprezentovat efekt toho, že entita hladoví. Pro knihovnu založenou na sparse set by jeden z typických přístupů byl vytvořit Starvation komponentu a přidat jí na každou entitu, která hladoví. Pro knihovny založené na arch type by ale tento přístup vedl ke ztrátě na výkonu. Jedním ze způsobů jak by se dal problém vyřešit je, že by každá entita co může hladovět obsahovala Hunger komponentu, ve které by byl flag Starvation který by rozhodoval o tom jestli daná entita hladoví.

Velký počet entit nám usnadní měření výkonu naší hry. Také nám zaručí, že systémy budou vykonávat netriviální množství práce. Zároveň si tím vyzkoušíme velký výkon, kterým ECS disponuje. Velký počet entit nám také implikuje větší herní svět.

Bude pro nás důležité aby naše hra nebyla příliš jednoduchá. Budeme se chtít vyhnout tomu, aby naše měření nebylo jenom dalším jednoduchým měřením výkonu. Ovšem taky si musíme dát pozor na to aby nebyla příliš složitá, jelikož to by odvádělo pozornost od problému, který řešíme.

% \\
% \xxx{Jaka charakteristiky by nase hra mela mit?}
% \\
% \xxx{- Velky pocet entit.}
% \\
% \xxx{- Hra musi byt dostatecne slozita aby vysledek mel prakticky smysl.}
% \\
% \xxx{- Hra nesmi byt prislis slozita aby slozitost neodvadela od problemu.}
% \\
% \xxx{- Hra by nemela vyuzivat casteho pridavani/odebirani komponent.}
% \\
% \xxx{-- Kdyby ano, nektere knihovny by byli zvyhodnene.}

\section{Simulace}
Žánrem naší hry bude simulace, ale řekněme si nejprve co to vlastně je. Simulace je žánr video her, který se snaží simulovat agenty v jejich prostředí. U nás agenti budou vesničané a jejich prostředí bude příroda.

Proč jsme zvolili právě simulaci? Pro simulaci je jednodušší (oproti ostatním žánrům) vymyslet herní design, který má velký počet entit a netriviálně velký svět. Další velkou výhodou je malí vliv uživatele na průběh hry, to nám umožní lépe provádět měření. Konkrétně v naší hře hráč bude spíše pozorovatel, který sleduje jak simulace probíhá.

% \\
% \xxx{Proc jsme zvolili simulaci?}
% \\
% \xxx{- Velky pocet entit.}
% \\
% \xxx{- Mali vliv uzivatele na prubeh hry.}
\chapter{Herní specifikace}
% intro
V této kapitole si představíme specifikaci naší hry. Cela hra se dá rozdělit na dva hlavní prvky. Konkrétně svět a vesnice. Specifikujeme si tedy každý z nich.

Prvně než si začneme specifikovat jednotlivé prvky, je nutné říci, že se jedná o 2D hru. Jak již bylo zmíněno, hráč v této hře zastává roli pozorovatele a celí svět pozoruje skrze kameru pohledem ze shora. Ovládá pouze pozici a přiblížení této kamery.

% svet
% - nahodne generovany
% - biomes
% - resources
% - - trees
% - - rocks
% - - iron
\section{Herní svět}
Herní svět je náhodně generovaný s pevnou velikostí a skládá se z několika biomů. Každý biom je reprezentován jinou barvou a nacházejí se v něm odlišné suroviny. Zároveň pro každý biom je určeno zda po něm můžou vesničané chodit nebo stavět. Herní svět obsahuje následující biomy:

\begin{enumerate}
    \item \textit{Voda} (\verb|Water|): Reprezentuje vodní plochu. Není na ní možné chodit ani stavět.
    \item \textit{Pláž} (\verb|Beach|): Jedná se o písečné oblasti v okolí \textit{vody}. Je na ní možné chodit i stavět.
    \item \textit{Louka} (\verb|Plain|): Na \textit{louce} je možné chodit i stavět. Je na nich možné najít \textit{jeleny}.
    \item \textit{Les} (\verb|Forest|): V \textit{lese} je možné chodit i stavět. Nachází se v něm velké množství \textit{stromů}.
    \item \textit{Hora} (\verb|Mountain|): Jedná se o kamenitou \textit{horu}. Je možné na ní chodit i stavět. Nachází se na ní \textit{naleziště kamene}.
    \item \textit{Vysoká hora} (\verb|High Mountain|): \textit{Vysoká hora je vždy obklopena \textit{horou}}. Je pokryta sněhovou plochou. Je možné na ní chodit ale není možné na ní stavět. Nachází se na ní \textit{naleziště železné rudy}.
\end{enumerate}

Jak již bylo zmíněno v jednotlivých biomech se můžou nacházet suroviny. Suroviny se rozdělují do několika typů. Typ suroviny udává jak dlouho trvá ji sklidit, jaký předmět lze získat po jejím sklizení a počet těchto předmětů, který lze získat po jejím sklizení. Hra obsahuje následující suroviny:

\begin{enumerate}
    \item \textit{Strom} (\verb|Tree|): Po pokácení \textit{stromu} lze získat předmět \textit{dřeva}. \textit{Stromy} se nacházejí ve velkém množství v biomu \textit{lesa}.
    \item \textit{Naleziště kamene} (\verb|Rock|): Sklizením \textit{naleziště kamene} je možné získat předmět \textit{kamene}. \textit{Naleziště kamene} se nacházejí v biomu \textit{hory}.
    \item \textit{Naleziště železné rudy} (\verb|Deposit|): Po sklizení \textit{naleziště železné rudy} je možné získat předmět \textit{železná ruda}. \textit{Naleziště železné rudy} se objevuje v biomu \textit{vysoká hora}.
    \item \textit{Jelen} (\verb|Deer|): \textit{Jeleni} na rozdíl od jiných typů surovin umí chodit. Náhodně se pohybují po biomu \textit{louky}, ve které se také objevují. Po sklizení jelena je možné získat předmět \textit{masa}.
\end{enumerate}

Předměty získané skrze sklizení suroviny, je dále možné zpracovat v příslušné budově. Každý předmět se může zpracovávat jinak dlouho. Hra obsahuje následující předměty.

\begin{enumerate}
    \item \textit{Dřevo} (\verb|Wood|): \textit{Dřevo} se dá dále zpracovat v \textit{prkno}.
    \item \textit{Kámen} (\verb|Stone|): \textit{Kámen} je možné přepracovat v \textit{cihlu}.
    \item \textit{Železná ruda} (\verb|Ore|): \textit{Železnou rudu} lze vypálit a tím získat \textit{železo}.
    \item \textit{Maso} (\verb|Meat|): \textit{Maso} je možné uvařit, tím lze získat \textit{jídlo}.
    \item \textit{Prkno} (\verb|Plank|): \textit{Prkna} je možné použít pro stavbu budov.
    \item \textit{Cihla} (\verb|Brick|): \textit{Cihlu} je také možné použít pro stavbu budov. Jedna se o lehce pokročilejší surovinu oproti \textit{prknu}.
    \item \textit{Železo} (\verb|Iron|): \textit{Železo} stejně jako \textit{Prkno} a \textit{Cihlu} lze také použít pro stavbu budov. Jedná se o pokročilou surovinu.
    \item \textit{Jídlo} (\verb|Food|): \textit{Jídlo} je předmět, který musí vesničané požívat, aby neumřeli hlady.
\end{enumerate}

% villages
% - villagers
% - houses
\section{Vesnice}
V herním světě se náhodně nacházejí vesnice. V každé se nachází několik budov a vesničané. Každá vesnice má přehled o předmětech kterými disponuje. Pro všechny vesnice je také pevně stanovená stavební fronta, kde pro každou budovu v této frontě je daná cena. Jakmile vesnice nasbírá dostatek surovin pro další budovu z této fronty, dojde k odečtení těchto surovin a následně se objeví příslušná budova. Budova se nemůže objevit moc blízko ani moc daleko od budov, které se již ve vesnici nacházejí.

Každý vesničan má svoje \textit{pracovní místo} (\verb|WorkPlace|). Vesničané vykonávají práci na základě své profese. Profese vesničanovi udává jaký předmět má pro vesnici obstarávat. Jeho práce spočívá v hledání a sklizení příslušné suroviny, kterou následně odnese do svého \textit{pracovního místa}, kde ji zpracuje a výsledný předmět odnese do \textit{skladiště}. Každý vesničan má právě jednu z následujících profesí:

\begin{enumerate}
    \item \textit{Dřevorubec} (\verb|Woodcutter|): Obstarává \textit{prkna}.
    \item \textit{Horník kamene} (\verb|StoneMiner|): Obstarává \textit{cihly}.
    \item \textit{Horník železa} (\verb|IronMiner|): Obstarává \textit{železo}.
    \item \textit{Lovec} (\verb|Hunter|): Obstarává \textit{jídlo}.
\end{enumerate}

Všechny vesnice začínají s \textit{hlavní budovou}, \textit{skladištěm} a \textit{budovou dřevorubce}. Budovy lze rozdělit do dvou kategorii a to speciální a budovy pro zpracování předmětů. V budovách pro zpracování předmětů je možné zpracovat předměty v jiné předměty, zároveň slouží jako \textit{pracovní místo} pro vesničany. Vždy když je postavena nová bude pro zpracování předmětů, objeví se nový vesničan, kterému bude přidělena příslušná profese podle této budovy. Speciální budovy slouží ke speciálním účelům a patří mezi ně všechny ostatní budovy. V každé vesnici se můžou nacházet následující typy budov:

\begin{enumerate}
    \item \textit{Hlavní budova} (\verb|MainBuilding|): Jedná se o speciální budovu, která reprezentuje danou vesnici. Každá vesnice má právě jednu hlavní budovu. Nemá žádný jiný zvláštní účel.
    \item \textit{Skladiště} (\verb|Stockpile|): Také se jedná o speciální budovu. Každá vesnice začíná s právě jedním skladištěm a nemůže jich mít více. Ve skladišti je možné skladovat jednotlivé předměty, kterými vesnice disponuje.
    \item \textit{Budova dřevorubce} (\verb|WoodcutterHut|): Budova pro zpracování předmětů, konkrétně \textit{prken}. Slouží jako pracovní místo pro \textit{dřevorubce}.
    \item \textit{Budova horníka} (\verb|MinerHut|): Budova pro zpracování předmětů, konkrétně \textit{cihel}. Slouží jako pracovní místo pro \textit{horníka kamene}.
    \item \textit{Kovárna} (\verb|Smithy|): Budova pro zpracování předmětů. konkrétně \textit{železa}. Slouží jako pracovní místo pro \textit{horníka železa}.
    \item \textit{Budova lovce} (\verb|HunterHut|): Budova pro zpracování předmětů, konkrétně \textit{jídla}. Slouží jako pracovní místo pro \textit{lovce}.
\end{enumerate}

\chapter{Hypotéza}

V této kapitole si rozebereme, jaké výsledky očekáváme od našeho měření.

Relativní výkon bude stejný jako u Ecs.CSharp.Benchmark.

Přiblížení cache?

Lépe vyjdou knihovny, které jsou založeny na archtype nebo sparse set.

Naopak knihovny, které dovolují pouze referenční datové typy pro komponenty vyjdou hůře.

% \\
% \xxx{Vysledek bude podobny jednoduchym testum z Ecs.CSharp.Benchmark.}
% \\
% \xxx{Vyjdou lepe ty knihovny co pouzivaji sparse set nebo archtype.}
% \\
% \xxx{Popsat co delaj spatne ty co vyjdou hure, napr pracuji s pointery na komponenty -> horsi vyuziti cache}
% \\
\chapter{Analýza přistupu k implementaci}

V této kapitole prozkoumáme přístupy k implementaci herních mechanik a rozeberm si proč jsme zvolili ty co jsme zvolili.

\section{Volba frameworku/enginu}

Prvně si řekněme co to vlastně je herní engine a herní framework.

Nýní si rozeberem populární herní enginy a frameworky pro C\#.

Proč jsme zvolili právě MonoGame?

% \\
% \xxx{Jake enginy/frameworky existujou pro C\#}
% \\
% \xxx{Proc jsme zvolili Monogame}
% \\

\section{Implementace AI}

Nyní si rozebereme populární přístupy pro implementaci AI ve hrách.

Nejpřímočarejší způsob je použít stavové automaty.

Pokročitejší zpusob jsou behavior trees.

Další pokročitelší způsob je Goal Oriented Action Planning (GOAP).

Proč jsme zvolili behavior trees?

Chování AI agentů může být popsané také přímo v systémech.

Ne vždy je zřejmé co by mělo být součástí AI.

% \\
% \xxx{jaka jsou mozsnosti?}
% \\
% \xxx{- behavior trees}
% \\
% \xxx{- state machines}
% \\
% \xxx{- GOAP}
% \\
% \xxx{proc jsme zvolili behavior trees?}
% \\
% \xxx{behavior primo v systemech a proc jsme to nepouzili}
% \\
% \xxx{Co by melo byt rizeno AI a co ne?}
% \\
% \xxx{- tezba surovin: pripsani do inventare vs damage and drop system}
% \\
% \xxx{- zpracovani surovin: pripsani do inventare vs resource processing system}
% \\

\section{ECS}

Některé systémy v ECS reagují pouze pokud je pro danou entitu splněna určitá podmínka.

Na první pohled by se mohlo zdát, že jsou tyto systémy špatné, jelikož podmínka je drahá operace.

Alternativní a pravděpodobně lepší způsob je použít eventy.

Rozeberme si jak by se s eventy v ECS pracovalo.

% \\
% \xxx{condition checking systems are ok}
% \\
% \xxx{eventy v ECS}
% \\

\section{Generovani terenu}

Rozeberme si dva přístupy pro generování terénu.

První způsob je generovat terén přímo na CPU.

Druhý způsob je pro generování terénu použít GPU.

% problem s path finding

Abychom si rozebrali řešení, bude nejprve nutné nejprve říci něco více o shaderech.

Compute shader můžem použít k řešení vžše zmíněného problému.

% \xxx{GPU vs CPU}

\section{Inline entity processor}

Inlinování funkcí je optimalizace, při které je volání funkce nahrazeno jejím obsahem.

% priklad, kdy dochazi k inlinovani

Kvůli měření je důležité, aby naše entity processory neměli ideálně žádný overhead.

Pro inlinování entity processorů, využíváma následujicí techniky.

% aggresive inlining, struktura + generika

% \\
% \xxx{Proc je to potreba}
% \\
% \xxx{Jak jsme toho dosahli}
% \\
\chapter{Popis implementace}
V této kapitole si blíže představíme implementaci naší hry.

\section{Abstrakce}
Prvně si rozeberem jak vypadá páteřní kód naší hry.

% Game, GameState, LevelState, Factory, LevelFactory

Nyní si přibližme naší abstrakční vrstvu Abstract.ECS.

% \xxx{Code base}
% \\
% \xxx{Abstract.ECS layer}

\section{Herní svět}
Herní svět je reprezentován třídou GameWorld.

O generování herního světa se stará třída GameWorldGenerator.

Rozeberme si co jsou to vlastně shadery.

Naše hra používá dva shadery.

% \xxx{GameWorld}
% \\
% \xxx{GameWorldGenerator}
% \\
% \xxx{Shaders}

\section{Herní data a jejich reprezentace}

Pro správu herních dat, slouží třída GlobalInstances.

% co jsou herni data, co trida dela

% \xxx{Data - biomes, items, and resources}

\section{Systémy a komponenty}

Nasledujicí seznam obsahuje a popisuje všechny herní systémy a komponenty.
\chapter{Analýza měření}
V této kapitole si rozebereme jaké existují přístupy k měření a na co při měření dbát ohled.

\section{Problematika měření}
Měření výkonu je komplexní proces, který by mohl být samostatnou prací sám o sobě.

Benchmark.NET je populární framework pro platformu .NET, který lze použít k měření výkonu.

\section{Přístupy k měření}
Nejprve je nutné si říct co to je vlastně game loop.

Nejpřímočarejší způsob je změřit jak dlouho trvá jedna iterace.

Lepší způsob by tedy bylo změřit n iterací a spočítat jejich průměr.

Rendering může být časově naročnější proces, rozeberme si jestli by měl či neměl být součastí měření.

Dalším způsobem jak měřit výkon, je změřit iterace jednotlivých systémů.
% proc to nedelame

Výkon je také možné měřit během hraní samotné hry.
% proc to nedelame

% \xxx{Delta time jedne iterace}
% \\
% \xxx{Delta time n iteracich}
% \\
% \xxx{Mel by byt rendering soucasti mereni}
% \\
% \xxx{Mereni prumerneho delta time jednotlivych ssytemu}
% \\
% \xxx{Mereni behem hrani}

\section{Na co při měření neberem ohled?}
Některé ECS knihovny pro získání lepšího výkonu používají paralelní queries.
% co to je
% nase hra je nepouziva protoze je z jednoduchosti navrzena single threaded

Další častá optimalizace, co některé ECS knihovny nabízejí jsou reaction systémy.
% co to je
% nepouzivame to protoze to nemaji vsechny knihovny

Jednotlivé ECS knihovny často používají další specifické výkonostní funkcionality.
% mi je nepouzivame protoze je nemaji vsechny knihovny

Další věc na kterou při měření neberem ohledy je přídávání a odebírání komponent.
% sparse set muzou by zvyhodnene
% nas herni design je prizpusoben tomu, ze to neovlivni mereni

% \xxx{Paralelni queries}
% \\
% \xxx{Reaction systemy}
% \\
% \xxx{Pridavani/odebirani komponent za runtime - design prizpusoben}
% \\
% \xxx{jine vykonostni vylepseni specificke pro danou ECS knihovnu}
\chapter{Experimenty a výsledky}
V této kapitole si prezentujem experimenty a jejich výsledky.

Schrneme si jak vlastně provádíme meření.

Následujicí přehled zobrazuje jednotlivé ECS knihovny a jejich výsledky.
% popsat prehled
% co vysledky znamenaji

\section{Analýza výsledků}
Nyní si stručně rozeberem jednotlivé ECS knihovny a stručně si řekneme proč se umístili tak jak se umístili.

% \xxx{Proc vysli takoveto vysledky}
% \\
% \xxx{Strucne rozebrat jednotlive knihovny a strucne popsat proc se umistili lepe/hure}
\chapter{Závěr}
Připoměnme si znovu, jaké byli vlastně cíle naší práce.

\section{Jak jsme splnili cíle práce?}
Cíl h1 se nám povedlo splnit.

Pomocí této hry jsme porovnali výkon populárních ECS knihoven pro C\# a tím splnili cíl h2.

Díky návrhu naší hry, jsme splnili v1 a prozkoumali jaké výhody ma navrhový vzor ECS.

V kapitole 2, jsme si rozebrali architektury, které je možné použít pro vývoj her a ukázali jsme si, jaké výhody má návrhový vzor ECS s porovnáním s návrhovým vzorem Komponent. Tím jsme splnili cíl v2.

\section{Co bychom mohli udělat lépe?}
Jedna z věcí co by mohla být lepší, je návrh samotné hry.
% kod + chyby

Hra by mohla být jednodušší.

\section{Možná budoucí vylepšení}
Hra by mohla být komplexnější.

Další možné vylepšení je podpora pro ECS knihovny zavislé na herním enginu Unity.

% \xxx{Co jsme chteli a ceho jsme dosahli}
% \\
% \xxx{cp bychom mohli udelat lepe}
% \\
% \xxx{future work}

\ifEN
\chapwithtoc{Bibliography}
\else
\chapwithtoc{Seznam použité literatury}
\fi

\printbibliography[heading=none]


%\appendix
%\chapter{Using CoolThesisSoftware}

Use this appendix to tell the readers (specifically the reviewer) how to use your software. A very reduced example follows; expand as necessary. Description of the program usage (e.g., how to process some example data) should be included as well.

To compile and run the software, you need dependencies XXX and YYY and a C compiler. On Debian-based Linux systems (such as Ubuntu), you may install these dependencies with APT:
\begin{Verbatim}
apt-get install \
  libsuperdependency-dev \
  libanotherdependency-dev \
  build-essential
\end{Verbatim}

To unpack and compile the software, proceed as follows:
\begin{Verbatim}
unzip coolsoft.zip
cd coolsoft
./configure
make
\end{Verbatim}

The program can be used as a C++ library, the simplest use is demonstrated in \cref{lst:ex}. A demonstration program that processes demonstration data is available in directory \verb|demo/|, you can run the program on a demonstration dataset as follows:
\begin{Verbatim}
cd demo/
./bin/cool_process_data data/demo1
\end{Verbatim}

After the program starts, control the data avenger with standard \verb-WSAD- controls.

\begin{listing}
\begin{lstlisting}
#include <CoolSoft.h>
#include <iostream>

int main() {
	int i;
	if(i = cool::ProcessAllData()) // returns 0 on error
		std::cout << i~<< std::endl;
	else
		std::cerr << "error!" << std::endl;
	return 0;
}
\end{lstlisting}
\caption{Example program.}
\label{lst:ex}
\end{listing}


% if your attachments are complicated, describe them in a separate appendix
%\chapwithtoc{Přílohy}
\label{chap:attachments}

\subsubsection{Přehled elektronických příloh}

\begin{enumerate}
    \item \texttt{src} -- Adresář obsahující zdrojové soubory a sestavený projekt měření. Součástí zdrojových souborů jsou dva projekty (WorldSimulator.csproj a WorldSimulator.Benchmarks.csproj).

    \item \texttt{build} -- Adresář obsahující sestavený projekt hry World Simulator.

    \item \texttt{README.md} -- Textový soubor obsahující opis seznamu příloh, dodatečné informace k seznamu příloh a kontakt na autora práce.
\end{enumerate}

\openright
\end{document}
